\subsection{Wrapping Lines}\label{sec:WrappingLines}
As we have an unlimited line length, it will not happen that an expression does not fit into a single line. Nevertheless you might consider to break long lines to improve the readability of the code. In this case, you should follow these general principles:

\subsubsection{Principles}
\begin{enumerate}[label=P\arabic*.]
\item{Only one statement per line!}
\item{Break after a comma.}
\item{Break before an operator.}
\item{Break before a dot.}
\item{Break after an opening parenthesis and before a closing one.}
\item{Prefer higher-level breaks to lower-level breaks.}
\item{Align the new line with the beginning of the expression at the same level on the previous line.}
\item{If the above rules lead to confusing code or to code that's squished up against the right margin, consider to break up the code into multiple parts, by introducing additional temporary variables.}
\end{enumerate}

\subsubsection*{Line Wrapping for Method Calls}
\keeplisting{4}
Here are some examples on how to break long method calls:
\begin{lstlisting}
//  USUAL
final var result = myMethod( longExpression1, longExpression2, 
    longExpression3, longExpression4, longExpression5 );
\end{lstlisting}
Using a single line for each argument makes especially sense in cases where one or more expressions are really long, or if a comment should be added to the argument:
\keeplisting{8}
\begin{lstlisting}
myProcedure
( 
    longExpression1,
    longExpression2,
    longExpression3,
    longExpression4, // some comment
    longExpression5
);
\end{lstlisting}

\keeplisting{9}
With a return value:
\begin{lstlisting}
final var result = myFunction
    ( 
        longExpression1,
        longExpression2,
        longExpression3,
        longExpression4, // some comment
        longExpression5
    );
\end{lstlisting}

\keeplisting{11}
When another method call is used as an argument\footnote{Usually, you should avoid to call another method in an arguments list, although sometimes this cannot be avoided (refer to \tqfullvref{sec:CodingConstructors} and there to section \tqfullref{sec:ConstructorHelpers}). But in most cases, introducing a local variable is the better option.}:
\keeplisting{9}
\begin{lstlisting}
final var result = myMethod( longExpression1,
    someMethod( longExpression2,
        longExpression3 ) );

final var result = myMethod1( longExpression1,
    someMethod( longExpression2,
        longExpression3 ),
    longExpression4,
    otherMethod( longExpression5 ) );
\end{lstlisting}

Next there are two examples of breaking an arithmetic expression. The first is preferred, since the break occurs outside the parenthesized expression, which is at a higher level.
\keeplisting{10}
\begin{lstlisting}
// PREFERRED
final var longName1 = longName2 * (longName3 + longName4 - longName5)
    + 4 * longname6;
           
longName1 = longName2 * (longName3 + longName4 - longName5)
    + 4 * longname6;
           
// AVOID!!!           
longName1 = longName2 * (longName3 + longName4
                       - longName5) + 4 * longname6;
\end{lstlisting}

\keeplisting{15}
Then we have a few examples for the indentation on long method declarations. First the conventional case:
\begin{lstlisting}
//  CONVENTIONAL INDENTATION
public final void myMethod( final int anArg, final Object anotherArg,
    final String yetAnotherArg, final Object andStillAnother )
{
    …
}   //  myMethod()

private final synchronized void horkingLongMethodName( final int arg,
   final Object anotherArg, final String yetAnotherArg, 
   final Object andStillAnotherArg )
{
    …
}   //  horkingLongMethodName()
\end{lstlisting}

Next, a \lstinline|throws| clause that will be either appended after the closing round bracket on the same line, or it is placed on its own line with an indentation of four blanks; in that case, the consecutive lines with parameters are indented by 8~spaces.
\keeplisting{13}
\begin{lstlisting}[numbers=left]
public final void myMethod( final int anArg, final Object anotherArg,
    final String yetAnotherArg, final Object andStillAnother,
    final Instant andTheFinalArg ) throws IllegalArgumentException
{
    …
}   //  myMethod()

// ADDITIONAL INDENTATION BECAUSE OF THE throws CLAUSE
public final void myMethod( final int anArg, final Object anotherArg,
        final String yetAnotherArg, final Object andStillAnother )
    throws IllegalArgumentException, IOException
{
    …
}   //  myMethod()
\end{lstlisting}
Usually, the second variant (starting at line~9) is better readable and should be preferred.

\keeplisting{11}
The final examples uses a single line for each argument:
\begin{lstlisting}[numbers=left]
private static final synchronized anotherHorkingLongMethodLine
(
    final int anArg,
    final Object anotherArg,  // a comment …
    final String yetAnotherArg,
    final Object andStillAnother
)
{
    …
}   //  anotherHorkingLongMethodLine()
\end{lstlisting}

\keeplisting{11}
Or again with a \lstinline|throws| clause:
\begin{lstlisting}[numbers=left]
private static synchronized anotherHorkingLongMethodLineThrows
(
    int anArg,
    Object anotherArg,
    @MyAnnotation String yetAnotherArg,
    Object andStillAnother
) throws IllegalArgumentException, IOException
{
    …
}   //  anotherHorkingLongMethodLineThrows()
\end{lstlisting}

For a method declaration, a comment for a formal parameter (as in line~4 of the first example) is usually obsolete, as the description for the parameters is a mandatory part of the documentation comments (see chapter \tqfullvref{sec:MethodComment} for the details).

But you should think about having each formal argument in a line of it own if you have annotations\index{annotation} on parameters, as shown in line~5 of the second example. Especially when the annotation\index{annotation} will have its own arguments.

\keeplisting{5}
Here are the acceptable ways to format ternary expressions\footnote{Read more about the ternary operator \lstinline|?| in chapter \tqfullvref{sec:TheTernaryOperator}.}:
\begin{lstlisting}[numbers=left]
final var alpha = (aLongBooleanExpression) ? beta : gamma;
                                           
final var alpha = (aLongBooleanExpression)
    ? beta
    : gamma;
\end{lstlisting}

And finally a bad example; the formula below calculates the date for Easter sunday for a given year, based on Gauss' Easter algorithm \autocite{WIKIPEDIA:DateOfEaster,WIKIPEDIA:Gaussche_Osterformel}.

\index{Gauss@Gauß, Carl Friedrich}\index{Lichtenberg, Heiner}\index{java.lang!IllegalArgumentException}\index{java.time!LocalDate}\index{java.time!LocalDate!minusDays()}\index{java.time!LocalDate!of()}\index{java.time!LocalDate!plusDays()}\index{java.time!Month!MARCH}\index{java.time!Year}\index{java.time!Year!getValue()}\index{java.time!Year!isAfter()}\index{java.time!Year!isBefore()}\index{org.tquadrat.foundation.lang!Objects}\index{org.tquadrat.foundation.lang!Objects!requireNonNullArgument()}
\begin{lstlisting}[numbers=left]
/**
 *  <p>{@summary This method calculates the date of Easter Sunday 
 *  for the given year.} The year has to be in the range from 1583 
 *  to 3900 (included).</p>
 *  <p>The resulting date is for the Gregorian calendar.</p>
 *  <p>The algorithm itself is not the original one published by Carl
 *  Friedrich Gauß first in 1800 (corrected version in 1816), but
 *  that one published by Heiner Lichtenberg in 1997.
 *
 *  @thanks Carl Friedrich Gauß
 *  @thanks Heiner Lichtenberg  
 *
 *  @param  year    The year for which the date of Easter Sunday is 
 *      wanted for.
 *  @return The date of Easter Sunday in the given year.
 *  
 *  @see <a href="https://de.wikipedia.org/wiki/Gau%C3%9Fsche_Osterformel">Gaußsche Osterformel</a>
 */
public static final LocalDate calcEasterDate( final Year year )
{
    //---* Check the arguments *-------------------------------------    
    if( requireNonNullArgument( year, "year" ).isBefore( Year.of( 1583 ) ) )
    {
        throw new IllegalArgumentException( "This method will work only for years greater than or equal to 1583" );
    }
    if( year.isAfter( Year.of( 3900 ) ) )
    {
        throw new IllegalArgumentException( "This method will work only for years less than or equal to 3900" );
    }

    /*
     * Initialise the return value with the last day of February and 
     * add the calculated number of days.
     */
    final var retValue = LocalDate.of( year.getValue(), MARCH, 1 )
        .minusDays( 1 )
        .plusDays( (21 + ((19 * (year.getValue() % 19) + (15 + (3 * (year.getValue() / 100) + 3) / 4 - (8 * (year.getValue() / 100) + 13) / 25)) % 30) - (((19 * (year.getValue() % 19) + (15 + (3 * (year.getValue() / 100) + 3) / 4 - (8 * (year.getValue() / 100) + 13) / 25)) % 30) + (year.getValue() % 19) / 11) / 29) + (7 - ((21 + ((19 * (year.getValue() % 19) + (15 + (3 * (year.getValue() / 100) + 3) / 4 - (8 * (year.getValue() / 100) + 13) / 25)) % 30) - (((19 * (year.getValue() % 19) + (15 + (3 * (year.getValue() / 100) + 3) / 4 - (8 * (year.getValue() / 100) + 13) / 25)) % 30) + (year.getValue() % 19) / 11) / 29) - (7 - (year.getValue() + year.getValue() / 4 + (2 - (3 * (year.getValue() / 100) + 3) % 4)) % 7)) % 7) );

    //---* Done *----------------------------------------------------
    return retValue;
}   //  calcEasterDate()
\end{lstlisting}

Even proper breaking of the term in line~37 will not help much to understand what's going on there. A much better formatted (and therefore more readable) version of the method \lstinline|calcEasterDate()| can be found in the listing \tqvref{listing:GaussEaster}, showing what is meant with \bdq{breaking up the code into multiple parts by introducing additional temporary variables}.

%==============================================================================
% End of file!!
%==============================================================================
